% Faithful transcription — original German. Transcribed from Vierteljahrsschrift
% der Naturforschenden Gesellschaft in Zürich, 15. Jahrgang, 3. Heft (1870), pp. 272–286:
% H. A. Schwarz, "Ueber einen Grenzübergang durch alternirendes Verfahren."
% Digital facsimile: Internet Archive / Biodiversity Heritage Library,
% archive.org/details/vierteljahrsschr15natu; printed pp. 272–286 correspond to leaves n279–n293.
%
% \origpage uses the printed page numbers (pp. 272–286). Running heads ("Schwarz,
% Grenzübergang durch alternirendes Verfahren.", "272", etc.) and the short closing rule
% on p. 286 are apparatus and not transcribed.
%
% TYPEFACE AND ORTHOGRAPHY. The text is set in Antiqua (Roman) type. Standard 19th-century
% German orthography is preserved verbatim: Ueber, alternirendes, convergirt, Constante,
% Theil, Werthe, dass, giebt, etc. Spaced type (Sperrdruck) used for proper names and key
% terms (Dirichlet, Riemann, Borchardt, Kronecker, Weber, Neumann, Weierstrass, Christoffel,
% Abel, zweistiefeligen, Parabel, Ellipse, Quadrat, Kreisbogendreieck, Nachtrag) is rendered
% as \emph{...}.
%
% FIGURE. The woodcut on p. 277 is embedded via \rmfigure{figures/fig-01.png} directly before
% the paragraph it illustrates (HOUSESTYLE R30).
%
% NOTATION. Laplace operator \Delta u = 0, partial derivatives \partial, centered dots in
% products preserved. Mathematical Fraktur \mathfrak{P}, \mathfrak{P}_1 preserved (HOUSESTYLE R13).
%
% APPARATUS. One compositor error on p. 286 ("Stèlle" with a grave accent) is flagged with \ednote.

\documentclass{article}
\usepackage{readmasters}
\begin{document}

\origpage{272}

\section*{Ueber einen Grenzübergang durch alternirendes Verfahren.}

\subsection*{Von \textbf{H.~A. Schwarz.}}

\emph{(Aus einem am 30.~Mai gehaltenen Vortrage.)}


Die unter dem Namen \emph{Dirichlet}'sches Princip bekannte Schlussweise, welche in gewissem Sinne als das Fundament des von \emph{Riemann} entwickelten Zweiges der Theorie der analytischen Funktionen angesehen werden muss, unterliegt, wie jetzt wohl allgemein zugestanden wird, hinsichtlich der Strenge sehr begründeten Einwendungen, deren vollständige Entfernung, soviel ich weiss, den Anstrengungen der Mathematiker bisher nicht gelungen ist.

Durch Fortsetzung einiger Untersuchungen, welche gewisse Arten von Abbildungsaufgaben betreffen, und von denen ein Theil im 70.~Band von \emph{Borchardt}'s Journal und in dem das Programm der eidgenössischen polytechnischen Schule für das Wintersemester 1869--70 begleitenden Aufsatz: „Zur Theorie der Abbildung“ veröffentlicht ist, bin ich auf ein Beweisverfahren geführt worden, durch welches, wie ich mich überzeugt zu haben glaube, alle Sätze, deren Beweis \emph{Riemann} in seinen veröffentlichten Abhandlungen mittelst des \emph{Dirichlet}'schen Princips zu führen gesucht hat, mit Strenge bewiesen werden können.

\origpage{273}

Die nachfolgende Mittheilung ist im Wesentlichen ein Auszug aus einer die Integration der partiellen Differentialgleichung $\Delta u = 0$ betreffenden Abhandlung, welche ich im November vorigen Jahres Hrn. \emph{Kronecker} und einigen andern mir nahestehenden Mathematikern mitgetheilt habe.

Es handelt sich wesentlich nur darum, den Nachweis der Existenz einer Funktion $u$ zu führen, welche für einen gewissen gegebenen Bereich $T$ der unabhängigen reellen Variablen $x$ und $y$ der partiellen Differentialgleichung
\[
\Delta u = \frac{\partial^2 u}{\partial x^2} + \frac{\partial^2 u}{\partial y^2} = 0
\]
und ausserdem gewissen vorgeschriebenen Grenz- und Unstetigkeitsbedingungen genügt.

Der Kürze wegen beschränke ich mich hier auf den Fall, in welchem die Nebenbedingungen nur Grenzbedingungen sind, in welchem also gefordert wird, es solle die Funktion $u$ stets endlich sein und längs der Begrenzung des Bereiches $T$ vorgeschriebene endliche Werthe haben, welche einer oder mehreren stetigen Folgen angehören. Auf diesen Fall kann nämlich durch das in der Folge mitzutheilende Verfahren der allgemeine Fall zurückgeführt werden.

Für die Anwendbarkeit des gedachten Beweisverfahrens ist es keineswegs erforderlich, die Voraussetzung zu machen, dass die Begrenzungslinie von $T$ nur eine endliche Anzahl von Ecken und im Allgemeinen in jedem Punkte einen endlichen bestimmten Krümmungsradius besitze, eine Voraussetzung, welche die HH. \emph{Weber} und \emph{Neumann} bei ihren auf dasselbe Ziel gerichteten Untersuchungen gemacht haben. (S. \emph{Borchardt}'s Journal, Band 71, pag.~29, und
\origpage{274}
Berichte der mathematisch-physischen Classe der K. Sächs. Gesellschaft der Wissenschaften, Sitzung vom 21. April 1870.) Es wird nicht einmal die Stetigkeit in der Aenderung der Richtung der Tangente der Begrenzungslinie erfordert; es genügt vielmehr zu wissen, dass die Begrenzungslinie sich in eine endliche Zahl von Stücken theilen lasse, so dass im Innern jedes dieser Stücke die Aenderung der Richtung der Tangente stets in demselben Sinne geschehe, wenn dieselbe auch unendlich oft sprungweise erfolgt, so dass also die Begrenzungslinie unendlich viele Ecken besitzen kann.

Auch Spitzen der Begrenzungslinie sind nicht ausgeschlossen. Für solche Spitzen, welche durch die Berührung zweier analytischen Linien entstehen, die in der Umgebung des Berührungspunktes den Charakter algebraischer Curven haben, habe ich die Untersuchung durchgeführt; um jedoch hier unnöthige Weitläufigkeiten zu vermeiden, ist im Folgenden auf das Auftreten von Spitzen keine Rücksicht genommen.

Das Gelingen des Beweises, dessen Grundgedanken hier mitgetheilt wird, beruht in letzter Instanz auf folgendem Hülfssatze:

Die Begrenzungslinie eines Bereiches $T$, für welchen es möglich ist, die partielle Differentialgleichung $\Delta u = 0$ beliebigen Grenzbedingungen gemäss zu integriren, werde in eine endliche Anzahl von Strecken (Theilen) getheilt. Diese mögen zu zwei Gruppen angeordnet werden, so dass in jeder Gruppe mindestens eine Strecke enthalten ist. Den einzelnen Strecken lege man, jenachdem sie der ersten oder zweiten Gruppe angehören, ungerade oder gerade Ordnungszahlen
\origpage{275}
bei und bezeichne die Punkte, welche die Strecken mit gerader Ordnungszahl von denen mit ungerader Ordnungszahl trennen, mit $P$. Im Innern von $T$ denke man sich eine endliche Anzahl analytischer Linien $L$ gegeben, welche mit den Strecken ungerader Ordnungszahl entweder keinen Punkt oder nur Endpunkte $P$ derselben gemein haben, ohne sie jedoch in diesen Punkten zu berühren.

Hierauf denke man sich für den Bereich $T$ eine Funktion $u$ bestimmt, welche der partiellen Differentialgleichung $\Delta u = 0$ genügt und in allen Punkten der Begrenzung von $T$ den Werth $0$ oder $+1$ hat, je nachdem die Ordnungszahl der Strecke, in deren Innerem der betreffende Punkt liegt, gerade oder ungerade ist. Dann ist die obere Grenze, beziehungsweise das Maximum aller Werthe, welche die Funktion $u$ längs der Linien $L$ annimmt, eine positive Zahl $q$, welche kleiner als $1$ ist.

Wird nun für denselben Bereich $T$ bei derselben Eintheilung der Begrenzung in Strecken mit gerader und ungerader Ordnungszahl und für dieselben Linien $L$ eine Funktion $u_1$ bestimmt, welche der Differentialgleichung $\Delta u_1 = 0$ genügt, auf der Begrenzung von $T$ längs der Strecken mit gerader Ordnungszahl den Werth Null hat, und deren längs der Strecken mit ungerader Ordnungszahl beliebig vorgeschriebener Werth dem absoluten Betrage nach die Grösse $g$ nicht überschreitet, so überschreitet der absolute Betrag der Werthe, welche die Funktion $u_1$ in den Punkten der Linien $L$ annehmen kann, nirgends den Werth $g \cdot q$, wo $q$ die im Vorhergehenden angegebene Bedeutung hat, also kleiner als $1$ ist. ---

\origpage{276}

Für die Fläche eines Kreises und für alle einfach zusammenhängenden Flächen, deren conforme Abbildung auf die Fläche eines Kreises bekannt ist, ist die Integration der partiellen Differentialgleichung $\Delta u = 0$ unter vorgeschriebenen Grenzbedingungen mit Schwierigkeit nicht verbunden. Hinsichtlich dieser Aufgabe möge es gestattet sein, auf einen kleinen in dieser Zeitschrift (S.~113--128 des laufenden Jahrganges) mitgetheilten Aufsatz Bezug zu nehmen; in demselben sind zwar Stetigkeitsunterbrechungen in der für die Funktion $u$ längs der Begrenzung vorgeschriebenen Werthenreihe der Kürze wegen ausdrücklich ausgeschlossen worden; indessen gelten die dort entwickelten Schlüsse \emph{mutatis mutandis} auch dann noch, wenn in einer endlichen Anzahl von Rand-Punkten die Reihe der vorgeschriebenen Randwerthe eine Unterbrechung der Stetigkeit erfährt.

Nachdem gezeigt ist, dass für eine Anzahl von einfacheren Bereichen die Differentialgleichung $\Delta u = 0$ beliebigen Grenzbedingungen gemäss integrirt werden kann, handelt es sich darum, den Nachweis zu führen, dass auch für einen weniger einfachen Bereich, der aus jenen auf gewisse Weise zusammengesetzt ist, die Integration der Differentialgleichung beliebigen Grenzbedingungen gemäss möglich ist. Zum Beweise dieses Satzes kann ein Grenzübergang dienen, welcher mit dem bekannten zur Herstellung eines luftverdünnten Raumes mittelst einer \emph{zweistiefeligen} Luftpumpe dienenden Verfahren grosse Analogie hat. Die Periode der Operation besteht nämlich in dem einen wie in dem andern Falle aus zwei alternirend zur Wirkung gelangenden Einzeloperationen, welche
\origpage{277}
zwar denselben Zweck haben, aber in Hinsicht auf die Art und Weise der Wirkung nicht identisch, sondern in gewissem Sinne symmetrisch sind.

Ein solcher Grenzübergang möge kurz Grenzübergang durch alternirendes Verfahren genannt werden.

\rmfigure{figures/fig-01.png}{}{Schematische Figur zweier sich überlappender Bereiche T_1 (Kreis) und T_2 (Quadrat) mit Schnittbereich T* und Randstücken L_0, L_1, L_2, L_3.}

Es seien gegeben zwei Bereiche $T_1$ und $T_2$, welche einen oder mehrere Bereiche $T^*$ gemeinsam haben und deren Begrenzungslinien sich nicht berühren. (In der beistehenden schematischen Figur ist $T_1$ die Fläche eines Kreises, $T_2$ die Fläche eines Quadrats.)

Das System aller Theile der Begrenzung von $T_1$, welche ausserhalb $T_2$ liegen, werde mit $L_0$, das System aller übrigen, innerhalb $T_2$ liegenden Theile mit $L_2$ bezeichnet.

Ebenso zerfällt die Begrenzung von $T_2$ in die Systeme $L_1$ und $L_3$, wenn nämlich mit $L_1$ das System aller Stücke, welche innerhalb des Gebietes $T_1$ liegen, mit $L_3$ das System aller Stücke, die ausserhalb $T_1$ liegen, bezeichnet wird.

Es wird vorausgesetzt, es sei sowohl für den Bereich $T_1$ als auch für den Bereich $T_2$ möglich, die Differentialgleichung $\Delta u = 0$ beliebigen Grenzbedingungen gemäss zu integriren; es handelt sich darum, zu zeigen, dass dies auch für den Bereich $T_1 + T_2 - T^* = T$ möglich ist, welcher die Bereiche $T_1$ und
\origpage{278}
$T_2$ als Theile enthält, bei welchem aber das den Gebieten $T_1$ und $T_2$ gemeinsame Gebiet $T^*$ nur einfach zu zählen ist.

Sowohl für das Gebiet $T_1$ und das System $L_1$ als auch für das Gebiet $T_2$ und das System $L_2$ sind die Bedingungen des vorher erwähnten Hülfssatzes erfüllt; im ersten Falle möge das System $L_0$, im zweiten das System $L_3$ an die Stelle der Gruppe der Strecken mit gerader Ordnungszahl treten. Es ist daher möglich, zwei Zahlen $q_1$ und $q_2$ zu bestimmen, welche die Rolle der Zahl $q$ in dem Hülfssatze vertreten und welche beide kleiner sind als $1$.

Dem Recipienten der Luftpumpe entspricht in Beibehaltung der obigen Analogie das Gebiet $T^*$, dem Innern der beiden Pumpencylinder entsprechen die Gebiete $T_1 - T^*$, $T_2 - T^*$, den Ventilen die Linien $L_1$ und $L_2$.

Es seien auf der Begrenzung von $T$, also längs $L_0$ und $L_3$, die Werthe für die Funktion $u$ willkürlich vorgeschrieben; $g$ sei die obere, $k$ sei die untere Grenze dieser Werthe; die Differenz $g - k$ werde mit $G$ bezeichnet.

Nun nehme man längs $L_2$ eine Werthenreihe willkürlich an, z.~B. in allen Punkten von $L_2$ den Werth $k$, und bestimme für das Gebiet $T_1$ eine Funktion $u_1$, welche längs $L_0$ die vorgeschriebenen Werthe, längs $L_2$ den Werth $k$ hat und im Innern von $T_1$ der Differentialgleichung $\Delta u_1 = 0$ genügt. Nach der über das Gebiet $T_1$ gemachten Voraussetzung gibt es eine solche Funktion. (Erster Zug des ersten Kolbens.)

Die Werthe, welche die Funktion $u_1$ längs $L_1$ hat, denke man sich fixirt und bestimme für das
\origpage{279}
Gebiet $T_2$ eine Funktion $u_2$, welche längs $L_3$ die vorgeschriebenen Werthe hat, längs $L_1$ mit der vorher bestimmten Funktion $u_1$ übereinstimmt und für welche $\Delta u_2 = 0$ ist. Nach der über das Gebiet $T_2$ gemachten Voraussetzung gibt es eine solche Funktion. (Erster Zug des zweiten Kolbens.)

Der Werth von $u_2 - u_1$ oder von $u_2 - k$ längs $L_2$ ist kleiner als $g - k = G$.

Man bestimme nun für das Gebiet $T_1$ eine Funktion $u_3$, welche längs $L_0$ die vorgeschriebenen Werthe hat, längs $L_2$ mit $u_2$ übereinstimmt und für welche $\Delta u_3 = 0$ ist. (Zweiter Zug des ersten Kolbens.)

Die Differenz $u_3 - u_1$ ist im Innern von $T_1$ in keinem Punkte negativ und dem absoluten Betrage nach kleiner als $G$, längs $L_1$ aber nach dem erwähnten Hülfssatze kleiner als $G \cdot q_1$, weil $u_3 - u_1$ längs $L_0$ den Werth Null hat und längs $L_2$ kleiner als $G$ ist.

Den Werth der Funktion $u_3$ längs $L_1$ denke man sich fixirt und für das Gebiet $T_2$ eine Funktion $u_4$ bestimmt, welche längs $L_1$ mit $u_3$ übereinstimmt, längs $L_3$ die vorgeschriebenen Werthe hat und für welche $\Delta u_4 = 0$ ist. (Zweiter Zug des zweiten Kolbens.)

Die Differenz $u_4 - u_2$ hat längs $L_3$ den Werth Null und ist längs $L_1$, wo sie mit $u_3 - u_1$ übereinstimmt, positiv und kleiner als $G \cdot q_1$; daher ist im Innern von $T_2$ $u_4 - u_2$ nirgends negativ und beständig kleiner als $G \cdot q_1$, längs $L_2$ aber kleiner als $G \cdot q_1 \cdot q_2$.

Durch Fortsetzung dieses alternirenden Verfahrens gelangt man zu einer Reihe von unendlich vielen Funktionen mit ungeradem und mit geradem Index. Die einen sind für das Gebiet $T_1$, die andern für das Gebiet $T_2$ so erklärt, dass sie beziehlich längs
\origpage{280}
$L_0$ und $L_3$ die vorgeschriebenen Werthe haben und im Innern der Gebiete, für welche sie erklärt sind, der partiellen Differentialgleichung $\Delta u = 0$ genügen.

Für das Gebiet $T^*$ sind sowohl die Funktionen mit ungeradem als die mit geradem Index erklärt und zwar stimmen dieselben abwechselnd längs $L_1$ und längs $L_2$ mit einander überein. Längs $L_1$ ist nämlich $u_{2n-1} = u_{2n}$ und längs $L_2$ $u_{2n+1} = u_{2n}$.

Es ist nun nicht schwer nachzuweisen, dass die Funktionen mit ungeradem und diejenigen mit geradem Index sich mit wachsendem Index bestimmten Grenzfunktionen $u'$ und $u''$ unbegrenzt nähern, welche durch die Gleichungen
\[
\begin{aligned}
u' &= u_1 + (u_3 - u_1) + (u_5 - u_3) + \dots \\
   &\qquad + (u_{2n+1} - u_{2n-1}) + \dots \text{ in inf.} \\
u'' &= u_2 + (u_4 - u_2) + (u_6 - u_4) + \dots \\
   &\qquad + (u_{2n+2} - u_{2n}) + \dots \text{ in inf.}
\end{aligned}
\]
erklärt sind. Die auf der rechten Seite stehenden Reihen convergiren unbedingt und für alle in Betracht kommenden Werthepaare $x, y$ in gleichem Grade; es ist nämlich
\[
\begin{aligned}
(u_{2n+1} - u_{2n-1}) & < G \cdot (q_1 \cdot q_2)^{n-1} \quad\text{und}\\
(u_{2n+2} - u_{2n}) & < G \cdot (q_1 \cdot q_2)^{n-1} \cdot q_1.
\end{aligned}
\]

Sowohl längs $L_1$ als längs $L_2$ ist $u' = u''$. Im Innern von $T_1$ ist $\Delta u' = 0$ im Innern von $T_2$ $\Delta u'' = 0$, daher ist für jeden Punkt von $T^*$ $u' = u''$, weil auf der ganzen Begrenzung von $T^*$ beide Funktionen mit einander übereinstimmen.

Es sind daher die beiden Funktionen $u'$ und $u''$ Werthe derselben Funktion $u$, welche für das ganze Gebiet $T = T_1 + T_2 - T^*$ erklärt ist, im Innern desselben der partiellen Differentialgleichung $\Delta u = 0$
\origpage{281}
genügt und auf der Begrenzung $L_0 + L_3$ die vorgeschriebenen Werthe annimmt.

Hiermit ist der Beweis für die oben ausgesprochene Behauptung angedeutet: unter den angegebenen Voraussetzungen ist es auch für den Bereich $T$ möglich, die partielle Differentialgleichung $\Delta u = 0$ willkürlich vorgeschriebenen Grenzbedingungen gemäss zu integriren. ---

Durch wiederholte Anwendung und geeignete Modifikation des erwähnten Grenzüberganges durch alternirendes Verfahren kann die Existenz einer Funktion $u$ für ein gegebenes Gebiet auch dann, wenn ausser den Grenzbedingungen noch Unstetigkeitsbedingungen, oder wie bei den \emph{Abel}'schen Integralen Unstetigkeitsbedingungen allein vorgeschrieben sind, in den Fällen dargethan werden, für welche \emph{Riemann} in seinen Abhandlungen die Existenz behauptet und mittelst des \emph{Dirichlet}'schen Princips zu beweisen gesucht hat.

Das erläuterte Beweisverfahren erstreckt sich nicht blos auf den Fall, in welchem die das Gebiet $T$ repräsentirende einfach oder mehrfach zusammenhängende \emph{Riemann}'sche Fläche in ihrer ganzen Ausdehnung in derselben Ebene oder auf derselben Kugelfläche enthalten ist, sondern gilt im Wesentlichen unverändert auch dann noch, wenn diese Fläche auf einer aus mehreren ebenen oder sphärischen Flächen zusammengesetzten Polyederoberfläche ausgebreitet ist.

Mit Hülfe dieser Erweiterung kann unter Anderem auch der Beweis geführt werden, dass ein einfach zusammenhängender auf einer solchen Polyederoberfläche ausgebreiteter Bereich conform abgebildet
\origpage{282}
werden kann auf die Fläche eines Kreises, wenn derselbe eine in sich zurückkehrende Begrenzungslinie besitzt, und auf die Fläche einer Kugel, wenn der Bereich ein einfach zusammenhängender und geschlossener Bereich ist.

Hiermit ist auch die Frage nach der Möglichkeit der Constantenbestimmung, auf welche die conforme Abbildung der einfach zusammenhängenden Oberfläche eines von ebenen Flächen begrenzten Polyeders auf eine Kugel zurückgeführt werden kann (s. \emph{Borchardt}'s Journal, Bd.~70, pag.~119), beantwortet.

Ein specieller Fall der soeben erwähnten Abbildungsaufgabe tritt ein, wenn es sich darum handelt, die einfach zusammenhängende Fläche eines von geradlinigen Strecken begrenzten ebenen Polygons auf die Fläche eines Kreises conform abzubilden, mag die Fläche des Polygons ganz im Endlichen liegen, oder den unendlich fernen Punkt der Ebene ein- oder mehrmals in ihrem Innern enthalten; auch Windungspunkte im Innern derselben sind nicht ausgeschlossen. Bei dieser Aufgabe liegt ebenfalls die einzige Schwierigkeit in dem Nachweise der Möglichkeit, eine gewisse Anzahl zum Theil reeller, zum Theil conjugirter complexer Constanten, von denen die abbildende Funktion abhängig gemacht wird, so zu bestimmen, dass allen Bedingungen der Aufgabe genügt wird.

Diese Schwierigkeit kann durch Anwendung eines von Herrn \emph{Weierstrass} entwickelten Verfahrens überwunden werden. Die Anwendung des obigen Grenzverfahrens bietet ein neues Mittel zur Ueberwindung derselben dar.

Aehnliches gilt von dem Nachweise der
\origpage{283}
Möglichkeit derjenigen Constantenbestimmung, auf welche die Aufgabe der conformen Abbildung einer von Kreisbogenstrecken begrenzten einfach zusammenhängenden Figur auf die Fläche eines Kreises zurückgeführt wird. (a.~a.~O. pag.~117.)

Auch in diesem Falle kann die Fläche in ihrem Innern Windungspunkte oder den unendlich fernen Punkt enthalten.

\emph{Nachtrag.} Vor Kurzem sind drei Abhandlungen des Herrn \emph{Christoffel} mir bekannt geworden (Annali di Matematica diretti da Brioschi e Cremona, Tomo IV${}^\circ$, pag.~1--9, Nachrichten von der K.~Gesellschaft der Wissenschaften zu Göttingen 1870, pag.~283--298 und 359--369), deren zu der obigen Mittheilung in Beziehung stehender Inhalt zu einigen Bemerkungen Anlass gibt.

Auf pag.~1 des vierten Bandes der Annali liest man: „\dots\ la determinazione delle temperature stazionarie sopra una superficie rettangolare $F$ non offre alcuna difficoltà, il problema associato delle temperature stazionarie sulla superficie complementare $F'$ è rimasto completamente inaccessibile ai metodi adoperati finquì\dots“

und auf pag.~284 der erwähnten Nachrichten: „\dots\ (man) gelangt dann zu einer merkwürdigen Gattung von Problemen, welche so auffallende Schwierigkeiten darbietet, dass die Lösung einer Aufgabe dieser Art ungeachtet aller Anstrengungen bisher nur in dem einzigen, durchaus elementaren Falle gelungen war, wo die Begrenzung von $\mathfrak{P}_1$ (--- die nach allen Richtungen in's Unendliche reichende Fläche, welche von einer Ebene übrig bleibt, wenn aus dieser
\origpage{284}
ein einfach zusammenhängendes, endliches Stück $\mathfrak{P}$ herausgeschnitten wird ---) ein Kreis ist. Namentlich sind, was im Folgenden seine Erklärung finden wird, alle Versuche gescheitert, irgend einen der besondern, vorzugsweise interessanten Fälle zu behandeln, wo $\mathfrak{P}_1$ von einer geradlinigten Figur begrenzt ist.“

Dieser Behauptung gegenüber erscheint es angemessen, auf einige einfache Beispiele aufmerksam zu machen, welche vielleicht auch an und für sich denen, welchen sie neu sind, einiges Interesse darbieten.

1.~Es sei in der Ebene der complexen Grösse $z$ eine \emph{Parabel} gegeben, deren Brennpunkt der Punkt $z = 0$, deren Scheitel der Punkt $z = +1$ ist. Durch die Funktion $Z = \operatorname{tg}^2\!\left(\frac{\pi}{4} \cdot \sqrt{z}\right)$ wird das Innere und durch die Funktion $Z_1 = \frac{2}{\sqrt{z}} - 1$ wird das Aeussere der Parabelfläche zusammenhängend und in den kleinsten Theilen ähnlich auf die Fläche je eines Kreises abgebildet, womit bekanntlich die betreffende Wärme-Aufgabe für das Aeussere und für das Innere der Parabel als gelöst zu betrachten ist.

2.~Es sei die Gleichung einer \emph{Ellipse} $\frac{x^2}{a^2} + \frac{y^2}{b^2} = 1$ gegeben; $a^2 - b^2 = 1$. Durch die Funktion $Z = \sin\mathrm{am}\left(\frac{2K}{\pi} \operatorname{arc\,sin} z\right)$, $q = \left(\frac{a - b}{a + b}\right)^2$, wird das Innere und durch die Funktion $Z_1 = \frac{z - \sqrt{z^2 - 1}}{a - b}$ wird das Aeussere der Ellipse auf die Fläche eines Kreises conform abgebildet, womit die betreffende Wärme-Aufgabe auch für diesen Fall als gelöst zu betrachten ist.

3.~Es sei gegeben ein \emph{Quadrat}, dessen Ecken
\origpage{285}
die vier Punkte $z = +1, +i, -1, -i$ bilden. Durch die Funktion $Z = \sin\mathrm{am}\, Kz$, $k = i$ wird das Innere des Quadrates auf die Fläche eines Kreises conform abgebildet und durch die Funktion
\[
z = -\frac{1}{C} \int^{Z_1} \frac{\sqrt{1 - Z_1^4}}{Z_1^2} \, dZ_1, \qquad C = \int_0^{\frac{\pi}{2}} \sqrt{\sin\varphi} \, d\varphi
\]
wird umgekehrt die Fläche des mit dem Radius $1$ um den Nullpunkt beschriebenen Kreises in der Ebene der complexen Grösse $Z_1$ auf das Aeussere jenes Quadrates conform abgebildet, wenn die complexe Integrationsvariable $Z_1$ auf jene Kreisfläche beschränkt wird und die Integrationsconstante aus der Bedingung bestimmt wird, dass $\lim \left(z - \frac{1}{CZ_1}\right)$ für $Z_1 = 0$ auch gleich $0$ sei. Durch die angegebenen Abbildungen darf aber die Wärme-Aufgabe für das Innere und für das Aeussere des Quadrates als gelöst angesehen werden. (Vrgl. \emph{Borchardt}'s Journ. Bd.~70, p.~115.)

4.~Es sei gegeben ein \emph{Kreisbogendreieck}. Die conforme Abbildung der Fläche eines Kreises auf das Innere und auf das Aeussere eines Kreisbogendreiecks ist mittelst hypergeometrischer Reihen ohne Schwierigkeit ausführbar. (Am angef. Orte p.~117.)

Die Zahl dieser Beispiele könnte noch vermehrt werden. ---

Es darf hier nicht unerwähnt gelassen werden, dass auch gegen den übrigen Inhalt der genannten Abhandlungen erhebliche Einwendungen geltend zu machen sind.

An einen vollständigen Beweis würde zweifelsohne die Forderung zu stellen sein, dass auch
\origpage{286}
bewiesen werde, es sei in der That möglich, sämmtliche Constanten, auf deren Bestimmung es ankommt, so zu bestimmen, dass allen Bedingungen der Aufgabe genügt wird.

Es ist ferner wohl zu beachten, dass in den beiden erwähnten Abhandlungen in den Göttinger Nachrichten die Untersuchung ausdrücklich auf den Fall beschränkt wird, in welchem die Begrenzungslinie des abzubildenden Bereiches durch eine „unzerfällbare Gleichung“ gegeben ist. Hierdurch werden also von vornherein sowohl die Fälle ausgeschlossen, in denen die Begrenzungslinie aus mehreren Stücken verschiedener analytischer Linien besteht, als auch der Fall, in welchem diese Linie an keiner Stèlle\ednote{So printed, with a grave accent on the first e: an apparent compositor's error for Stelle.} den Charakter einer algebraischen Curve besitzt.

Endlich ist nicht zu übersehen, dass in einer sehr grossen Anzahl von Fällen, zu welchen auch der einfache Fall der conformen Abbildung des Innern einer \emph{Ellipse} auf das Innere eines Kreises gehört, die in den genannten Abhandlungen aufgestellten Schlussformeln mit einer ungehobenen Schwierigkeit noch behaftet bleiben, ein Umstand, welcher es als fraglich erscheinen lässt, ob diese Formeln als ein Beweis für die Möglichkeit der Lösung der gestellten allgemeinen Aufgabe angesehen werden dürfen.

August 1870.

\end{document}
